\expandafter\gdef\csname SNside1\endcsname{\node[anchor=north west,inner sep=0pt,text width=48mm,font=\SideFont,align=left] (SNsidebox) at ([xshift=138mm,yshift=-406.28pt]current page.north west) {\textcolor{SNink}{\textbf{课程连接 · 分层连接操作系统}}\par\vspace{1.5mm}Cache利用短期访问局部性，虚拟存储还承担地址空间、隔离和换页。两者都能体现层次思想，但缓存对象与缺失处理不同：Cache常缓存数据块，TLB缓存映射，页表记录虚页到页框关系。};\pgfextracty{\SNsideheight}{\pgfpointdiff{\pgfpointanchor{SNsidebox}{south}}{\pgfpointanchor{SNsidebox}{north}}}\typeout{SNHEIGHT:1:0:\the\SNsideheight}}
\expandafter\gdef\csname SNside2\endcsname{\node[anchor=north west,inner sep=0pt,text width=48mm,font=\SideFont,align=left] (SNsidebox) at ([xshift=138mm,yshift=-95.00pt]current page.north west) {\textcolor{SNink}{\textbf{最小反例 · 动态不等于访问变化}}\par\vspace{1.5mm}最小反例：一个程序只反复读同一DRAM单元，芯片仍需刷新。动态描述电荷保持机制，而非程序写入频度；随机描述按地址选取能力，也不意味着无需刷新。把这两条分类轴分开。};\pgfextracty{\SNsideheight}{\pgfpointdiff{\pgfpointanchor{SNsidebox}{south}}{\pgfpointanchor{SNsidebox}{north}}}\typeout{SNHEIGHT:2:1:\the\SNsideheight}}
\expandafter\gdef\csname SNside3\endcsname{\node[anchor=north west,inner sep=0pt,text width=48mm,font=\SideFont,align=left] (SNsidebox) at ([xshift=138mm,yshift=-95.00pt]current page.north west) {\textcolor{SNink}{\textbf{最小反例 · 已安装容量不决定MAR}}\par\vspace{1.5mm}最小反例：机器给出64MiB字节地址空间，只安装32MiB内存。MAR仍需容纳26位地址，而非按已安装容量只取25位。芯片单元地址、可编址空间、已安装容量是三个问题。};\pgfextracty{\SNsideheight}{\pgfpointdiff{\pgfpointanchor{SNsidebox}{south}}{\pgfpointanchor{SNsidebox}{north}}}\typeout{SNHEIGHT:3:3:\the\SNsideheight}}
\expandafter\gdef\csname SNside4\endcsname{\node[anchor=north west,inner sep=0pt,text width=48mm,font=\SideFont,align=left] (SNsidebox) at ([xshift=138mm,yshift=-95.00pt]current page.north west) {\textcolor{SNink}{\textbf{不变量 · 位扩展共享片选}}\par\vspace{1.5mm}同一地址上的多片芯片分担不同数据位，因此片选和片内地址应同步；不同地址组才由高位译码区分。先画组，再连数据线，能避免把位扩展误画成容量分区。};\pgfextracty{\SNsideheight}{\pgfpointdiff{\pgfpointanchor{SNsidebox}{south}}{\pgfpointanchor{SNsidebox}{north}}}\typeout{SNHEIGHT:4:5:\the\SNsideheight}}
\expandafter\gdef\csname SNside5\endcsname{\node[anchor=north west,inner sep=0pt,text width=48mm,font=\SideFont,align=left] (SNsidebox) at ([xshift=138mm,yshift=-95.00pt]current page.north west) {\textcolor{SNink}{\textbf{错项排除 · 平均旋转等半圈}}\par\vspace{1.5mm}转速r转每分意味着一圈60/r秒；均匀随机角位置的平均等待为半圈，传一扇区另按一圈中该扇区比例计算。SSD没有这段机械旋转过程，不能给它机械套用半圈等待。};\pgfextracty{\SNsideheight}{\pgfpointdiff{\pgfpointanchor{SNsidebox}{south}}{\pgfpointanchor{SNsidebox}{north}}}\typeout{SNHEIGHT:5:6:\the\SNsideheight}}
\expandafter\gdef\csname SNside6\endcsname{\node[anchor=north west,inner sep=0pt,text width=48mm,font=\SideFont,align=left] (SNsidebox) at ([xshift=138mm,yshift=-95.00pt]current page.north west) {\textcolor{SNink}{\textbf{最小反例 · 循环轮数不是访问次数}}\par\vspace{1.5mm}最小反例：一条a{[}k{]}=a{[}k{]}+32每轮读一次又写一次，因此四元素八次数据访问。空Cache顺序扫描时一块首次读缺失，后续写命中。分母错用四轮，就把正确12.5\%误算成25\%。};\pgfextracty{\SNsideheight}{\pgfpointdiff{\pgfpointanchor{SNsidebox}{south}}{\pgfpointanchor{SNsidebox}{north}}}\typeout{SNHEIGHT:6:8:\the\SNsideheight}}
\expandafter\gdef\csname SNside7\endcsname{\node[anchor=north west,inner sep=0pt,text width=48mm,font=\SideFont,align=left] (SNsidebox) at ([xshift=138mm,yshift=-219.44pt]current page.north west) {\textcolor{SNink}{\textbf{最小反例 · 命中会改变LRU}}\par\vspace{1.5mm}最小反例：两路一组访问A、B、A、C。LRU在第三次命中A后淘汰B，FIFO仍按装入次序淘汰A。只记录装入顺序无法模拟LRU，命中也必须更新最近使用状态。};\pgfextracty{\SNsideheight}{\pgfpointdiff{\pgfpointanchor{SNsidebox}{south}}{\pgfpointanchor{SNsidebox}{north}}}\typeout{SNHEIGHT:7:10:\the\SNsideheight}}
\expandafter\gdef\csname SNside8\endcsname{\node[anchor=north west,inner sep=0pt,text width=48mm,font=\SideFont,align=left] (SNsidebox) at ([xshift=138mm,yshift=-95.00pt]current page.north west) {\textcolor{SNink}{\textbf{最小反例 · 映射没缓存不等于没驻留}}\par\vspace{1.5mm}最小反例：TLB未命中，但页表存在位为1。查表后即可得到页框号，无需从磁盘调页。TLB缺失只说明映射未缓存，缺页才说明合法目标页当前不驻留。\par\smallskip\begin{tikzpicture}[x=1mm,y=1mm,>=Stealth,node distance=5mm]
\node[draw=SNink!50,fill=SNblue,rounded corners=1mm,text width=34mm,align=center,inner sep=1.2mm,font=\fontsize{7.5}{9}\selectfont] (m0) at (0,0) {TLB缺失};
\node[draw=SNink!50,fill=SNblue,rounded corners=1mm,text width=34mm,align=center,inner sep=1.2mm,font=\fontsize{7.5}{9}\selectfont] (m1) at (0,-11) {页已驻留};
\node[draw=SNink!50,fill=SNblue,rounded corners=1mm,text width=34mm,align=center,inner sep=1.2mm,font=\fontsize{7.5}{9}\selectfont] (m2) at (0,-22) {查表继续};
\draw[->,SNink] (m0) -- (m1);
\draw[->,SNink] (m1) -- (m2);
\end{tikzpicture}
\par{\fontsize{7.5}{9}\selectfont TLB未命中且存在位为1，不发生缺页。}};\pgfextracty{\SNsideheight}{\pgfpointdiff{\pgfpointanchor{SNsidebox}{south}}{\pgfpointanchor{SNsidebox}{north}}}\typeout{SNHEIGHT:8:14:\the\SNsideheight}}
\expandafter\gdef\csname SNside9\endcsname{\node[anchor=north west,inner sep=0pt,text width=48mm,font=\SideFont,align=left] (SNsidebox) at ([xshift=138mm,yshift=-95.00pt]current page.north west) {\textcolor{SNink}{\textbf{不变量 · 写回与写分配独立}}\par\vspace{1.5mm}直写或回写回答命中时怎样更新下一层，分配或不分配回答缺失时是否调入块。常见搭配不是逻辑上的唯一搭配。读清题设组合，再决定缺失数、主存写次数和脏块处理。};\pgfextracty{\SNsideheight}{\pgfpointdiff{\pgfpointanchor{SNsidebox}{south}}{\pgfpointanchor{SNsidebox}{north}}}\typeout{SNHEIGHT:9:11:\the\SNsideheight}}
\expandafter\gdef\csname SNside10\endcsname{\node[anchor=north west,inner sep=0pt,text width=48mm,font=\SideFont,align=left] (SNsidebox) at ([xshift=138mm,yshift=-95.00pt]current page.north west) {\textcolor{SNink}{\textbf{不变量 · 段界限先约束偏移}}\par\vspace{1.5mm}段式地址包含段号与段内偏移，应先核段表和偏移是否在界限内，再形成物理地址或进入分页。扩大基址不能修复越界偏移；非法地址与合法页未驻留也不能混成同一种缺页。};\pgfextracty{\SNsideheight}{\pgfpointdiff{\pgfpointanchor{SNsidebox}{south}}{\pgfpointanchor{SNsidebox}{north}}}\typeout{SNHEIGHT:10:15:\the\SNsideheight}}
\expandafter\gdef\csname SNside11\endcsname{\node[anchor=north west,inner sep=0pt,text width=48mm,font=\SideFont,align=left] (SNsidebox) at ([xshift=138mm,yshift=-95.00pt]current page.north west) {\textcolor{SNink}{\textbf{错项排除 · 缺页后停止本轮查数据}}\par\vspace{1.5mm}若页表判断目标合法但未驻留，本轮尚无可用页框号，不能拿旧字段继续查数据Cache。先完成缺页处理、更新映射并重执行访问。题目未给主存字节时，也不能凭路径猜出数据值。};\pgfextracty{\SNsideheight}{\pgfpointdiff{\pgfpointanchor{SNsidebox}{south}}{\pgfpointanchor{SNsidebox}{north}}}\typeout{SNHEIGHT:11:16:\the\SNsideheight}}
\expandafter\gdef\csname SNside12\endcsname{\node[anchor=north west,inner sep=0pt,text width=48mm,font=\SideFont,align=left] (SNsidebox) at ([xshift=138mm,yshift=-95.00pt]current page.north west) {\textcolor{SNink}{\textbf{错项排除 · 缺页后停止本轮查数据}}\par\vspace{1.5mm}若页表判断目标合法但未驻留，本轮尚无可用页框号，不能拿旧字段继续查数据Cache。先完成缺页处理、更新映射并重执行访问。题目未给主存字节时，也不能凭路径猜出数据值。};\pgfextracty{\SNsideheight}{\pgfpointdiff{\pgfpointanchor{SNsidebox}{south}}{\pgfpointanchor{SNsidebox}{north}}}\typeout{SNHEIGHT:12:16:\the\SNsideheight}}
\expandafter\gdef\csname SNside13\endcsname{\node[anchor=north west,inner sep=0pt,text width=48mm,font=\SideFont,align=left] (SNsidebox) at ([xshift=138mm,yshift=-95.00pt]current page.north west) {\textcolor{SNink}{\textbf{错项排除 · 缺页后停止本轮查数据}}\par\vspace{1.5mm}若页表判断目标合法但未驻留，本轮尚无可用页框号，不能拿旧字段继续查数据Cache。先完成缺页处理、更新映射并重执行访问。题目未给主存字节时，也不能凭路径猜出数据值。};\pgfextracty{\SNsideheight}{\pgfpointdiff{\pgfpointanchor{SNsidebox}{south}}{\pgfpointanchor{SNsidebox}{north}}}\typeout{SNHEIGHT:13:16:\the\SNsideheight}}
\expandafter\gdef\csname SNside14\endcsname{\node[anchor=north west,inner sep=0pt,text width=48mm,font=\SideFont,align=left] (SNsidebox) at ([xshift=138mm,yshift=-95.00pt]current page.north west) {\textcolor{SNink}{\textbf{错项排除 · 缺页后停止本轮查数据}}\par\vspace{1.5mm}若页表判断目标合法但未驻留，本轮尚无可用页框号，不能拿旧字段继续查数据Cache。先完成缺页处理、更新映射并重执行访问。题目未给主存字节时，也不能凭路径猜出数据值。};\pgfextracty{\SNsideheight}{\pgfpointdiff{\pgfpointanchor{SNsidebox}{south}}{\pgfpointanchor{SNsidebox}{north}}}\typeout{SNHEIGHT:14:16:\the\SNsideheight}}
\expandafter\gdef\csname SNside15\endcsname{\node[anchor=north west,inner sep=0pt,text width=48mm,font=\SideFont,align=left] (SNsidebox) at ([xshift=138mm,yshift=-296.51pt]current page.north west) {\textcolor{SNink}{\textbf{错项排除 · 缺页后停止本轮查数据}}\par\vspace{1.5mm}若页表判断目标合法但未驻留，本轮尚无可用页框号，不能拿旧字段继续查数据Cache。先完成缺页处理、更新映射并重执行访问。题目未给主存字节时，也不能凭路径猜出数据值。};\pgfextracty{\SNsideheight}{\pgfpointdiff{\pgfpointanchor{SNsidebox}{south}}{\pgfpointanchor{SNsidebox}{north}}}\typeout{SNHEIGHT:15:16:\the\SNsideheight}}
\expandafter\gdef\csname SNside16\endcsname{\node[anchor=north west,inner sep=0pt,text width=48mm,font=\SideFont,align=left] (SNsidebox) at ([xshift=138mm,yshift=-95.00pt]current page.north west) {\textcolor{SNink}{\textbf{课程连接 · 从路径列事件}}\par\vspace{1.5mm}先列访问哪些层级、命中或缺失时发生哪些访问，再按模型计算平均时间；不要把不同路径的额外时间重复计入。};\pgfextracty{\SNsideheight}{\pgfpointdiff{\pgfpointanchor{SNsidebox}{south}}{\pgfpointanchor{SNsidebox}{north}}}\typeout{SNHEIGHT:16:18:\the\SNsideheight}}
\expandafter\gdef\csname SNside17\endcsname{\node[anchor=north west,inner sep=0pt,text width=48mm,font=\SideFont,align=left] (SNsidebox) at ([xshift=138mm,yshift=-95.00pt]current page.north west) {\textcolor{SNink}{\textbf{不变量 · 相联程度不直接等于延迟}}\par\vspace{1.5mm}全相联能并行比较多行，组相联只比较一个组；实际延迟还受比较器、布线和编码影响。固定实现才可比较快慢。映射限制候选位置，替换选择牺牲行，电路时序决定访问时间。};\pgfextracty{\SNsideheight}{\pgfpointdiff{\pgfpointanchor{SNsidebox}{south}}{\pgfpointanchor{SNsidebox}{north}}}\typeout{SNHEIGHT:17:17:\the\SNsideheight}}
\expandafter\gdef\csname SNside18\endcsname{\node[anchor=north west,inner sep=0pt,text width=48mm,font=\SideFont,align=left] (SNsidebox) at ([xshift=138mm,yshift=-95.00pt]current page.north west) {\textcolor{SNink}{\textbf{课程连接 · 页与擦除块}}\par\vspace{1.5mm}闪存按页读取或编程，按块擦除；逻辑覆盖常转为异地写入并更新映射，磨损均衡不能消除擦除次数上限。};\pgfextracty{\SNsideheight}{\pgfpointdiff{\pgfpointanchor{SNsidebox}{south}}{\pgfpointanchor{SNsidebox}{north}}}\typeout{SNHEIGHT:18:20:\the\SNsideheight}}
\expandafter\gdef\csname SNside19\endcsname{\node[anchor=north west,inner sep=0pt,text width=48mm,font=\SideFont,align=left] (SNsidebox) at ([xshift=138mm,yshift=-95.00pt]current page.north west) {\textcolor{SNink}{\textbf{课程连接 · 位扩展与字扩展}}\par\vspace{1.5mm}位扩展让多片提供同一地址的不同数据位；字扩展选择不同芯片提供更多地址，两者可组合。};\pgfextracty{\SNsideheight}{\pgfpointdiff{\pgfpointanchor{SNsidebox}{south}}{\pgfpointanchor{SNsidebox}{north}}}\typeout{SNHEIGHT:19:21:\the\SNsideheight}}
\expandafter\gdef\csname SNside20\endcsname{\node[anchor=north west,inner sep=0pt,text width=48mm,font=\SideFont,align=left] (SNsidebox) at ([xshift=138mm,yshift=-95.00pt]current page.north west) {\textcolor{SNink}{\textbf{课程连接 · 命中也会改变LRU}}\par\vspace{1.5mm}LRU命中后需把命中块更新为最近使用；FIFO命中一般不改变装入顺序。每次访问后写出状态，避免只数缺失。};\pgfextracty{\SNsideheight}{\pgfpointdiff{\pgfpointanchor{SNsidebox}{south}}{\pgfpointanchor{SNsidebox}{north}}}\typeout{SNHEIGHT:20:23:\the\SNsideheight}}
\expandafter\gdef\csname SNside21\endcsname{\node[anchor=north west,inner sep=0pt,text width=48mm,font=\SideFont,align=left] (SNsidebox) at ([xshift=138mm,yshift=-95.00pt]current page.north west) {\textcolor{SNink}{\textbf{课程连接 · 地址划分顺序}}\par\vspace{1.5mm}依题意先确定块内偏移，再根据组数确定组号；比较标记前还要验证有效位。};\pgfextracty{\SNsideheight}{\pgfpointdiff{\pgfpointanchor{SNsidebox}{south}}{\pgfpointanchor{SNsidebox}{north}}}\typeout{SNHEIGHT:21:26:\the\SNsideheight}}
\expandafter\gdef\csname SNside22\endcsname{\node[anchor=north west,inner sep=0pt,text width=48mm,font=\SideFont,align=left] (SNsidebox) at ([xshift=138mm,yshift=-95.00pt]current page.north west) {\textcolor{SNink}{\textbf{课程连接 · 地址划分顺序}}\par\vspace{1.5mm}依题意先确定块内偏移，再根据组数确定组号；比较标记前还要验证有效位。};\pgfextracty{\SNsideheight}{\pgfpointdiff{\pgfpointanchor{SNsidebox}{south}}{\pgfpointanchor{SNsidebox}{north}}}\typeout{SNHEIGHT:22:26:\the\SNsideheight}}
\AddToShipoutPictureFG{\begin{tikzpicture}[remember picture,overlay]\ifcsname SNside\thepage\endcsname\csname SNside\thepage\endcsname\fi\end{tikzpicture}}